\section{The original comparison packet and its transport to one fitting realization}
\label{sec:comparison-packet}

This chapter constructs the packet used immediately before the actual readers of
RDR1--RDR22.  The formal source is fixed at commit
\texttt{7cbf7a0c32fca1badf05931879a37c3f8a2b03b4}; its correctness is a standing
assumption.  What follows is an independent mathematical reconstruction of the
constructors, with their original inputs retained.  In particular, a supplied
geometric realization is never replaced by an unrelated isomorphic model.
The Hilbert spaces that are introduced below are coordinates for explicitly
specified original subspaces, with maps back to those subspaces included in
every assertion.

The principal conclusions are \texttt{exists\_pkT}, the full
\texttt{rowPk\_nonempty}, and \texttt{exists\_moved}.  They are conditional on
the same original \texttt{Inputs}, \texttt{Pins}, \texttt{Holds}, and, for the
last conclusion, \texttt{ExAt}, \texttt{InGeo}, and \texttt{Fits}.
They do not prove existence of those data.  Nor do they identify an arbitrary
later congruence parameter with the parameter extracted from the realization.
The latter equality remains exactly the guard in RDR22.  No additional
row-integrability hypothesis is introduced: the original row integrals already
have the supplier RPI14, and RPI15 supplies the first division comparison.

\subsection{Fixed objects and their complex coordinates}

% BEGIN PKT1
\paragraph{PKT1 (the original interfaces and selection order).}\label{claim:PKT1}
Fix the original $D$, finite group $G$, three-stage family $A_3$, coefficient
field $K=D.\mathrm{what}_2.\mathrm{LL}$ with its star-preserving embedding $\sigma$, intermediate
number field $K_s$, and original records
\[
 R:\mathrm{RealDat3}(D,A_3,K_s),\qquad F:\mathrm{Face3}(D,A_3).
\]
Assume the original $\mathrm{Inputs}(D)$, $\mathrm{Pins}(R,F,c_G)$ and
$\mathrm{Holds}(R,F)$, and that the image of the original coefficient ring
$O$ in $K$ lies in $K_s$.  Its elements are integral by the original $D.O_{\rm int}$ field.  Fix $y:\mathrm{Hd}(R)$ and, before constructing
any comparison maps, fix matrices $m_A(b,\mu)$ and $m_L(\mu,j)$ satisfying
all original \texttt{Mat1.Laws}.  Assume the literal equalities
\[
\begin{split}
 m_A(b,y.\mathrm{tup}(t))&=\mathrm{rMat}(A_3.A,P_b,\psi_b,t),\\
 m_L(y.\mathrm{tup}(t),i+1)&=\mathrm{lMat}(A_3.A,P_0,\psi_0,i,t).
\end{split}\tag{PKT.1}
\]
Here $b=\mathrm{true}$ means $(P_1,\psi_1)$ and $b=\mathrm{false}$ means
$(P_0,\psi_0)$.  Write $n_b=\mathrm{nP}(F,b)=|P_b|>0$.

The packet to be constructed has one finite type catalogue, all complexified
spaces $s_A,s_L$, both row chains, one lift family, the full multiplicity
integer $D_{\rm all}$, the real number $\delta_p$, and the sharp original
row coordinates.  Its accompanying \texttt{PkT} contains the catalogue
certificate, every complexification tie, both chain comparisons, the full
lift comparison, and the equality defining $D_{\rm all}$.  The equality
$p.m_L=m_L$ is literal.  Thus none of these choices may be made again after
forming a reader.  These are precisely the fields in the source definitions;
this paragraph fixes the scope of all subsequent constructions.
% END PKT1

% BEGIN PKT2
\paragraph{PKT2 (the exact complexification and lattice).}\label{claim:PKT2}
Let $m=(n,H,E,N,\mathrm{op})$ satisfy the matrix laws at $\sigma$.
Choose an invertible complex matrix $S$ with $S^*S=\sigma(H)$ and put
\[
 j(x)=S\sigma(x),\qquad \rho(B)=S\sigma(B)S^{-1}.
 \tag{PKT.2}
\]
The spectral theorem, or Gram--Schmidt for the positive Hermitian form, gives
$S$; when $n=0$ these are the unique empty matrices.  The images of the
coordinate vectors are a complex basis, and
$\langle jx,jz\rangle=\sigma(x^*Hz)$.  The matrix law $E=\mathrm{op}(e)$,
central multiplicativity and $e^2=e$ imply $E^2=E$.  The stated selfadjoint
law gives symmetry of $\rho(E)$.  Thus $e_m=\rho(E)$ is an orthogonal
idempotent and $j(Ex)=e_mj(x)$.

The full $O_K$-lattice $N$ is finitely generated over $\mathbb Z$: choose
$O_K$-generators and multiply them by an integral basis of $O_K$.  Let
$u_1,\ldots,u_a$ be the resulting generators.  Define the required
$O$-lattice by the finite generating set $\{j(u_r)\}$ and carrier
\[
 \Lambda=j(N).\tag{PKT.3}
\]
This is exactly the $O$-span of those generators.  Indeed every coefficient
coming from $O$ is integral in $K$, hence lies in $O_K$, so the $O$-span is
contained in $j(N)$; conversely integer coefficients already generate $N$.
Since $KN=K^n$, these vectors span the whole complex space after applying
$j$ and scalar extension.  This proves every generator, carrier and spanning
clause of \texttt{OLatTie}, and also stability under every $\sigma(a)$,
$a\in O_K$.  It uses no surjectivity of $O\to O_K$.
% END PKT2

% BEGIN PKT3
\paragraph{PKT3 (the actual central group-algebra action).}\label{claim:PKT3}
Only central multiplicativity is assumed for $\mathrm{op}$; we therefore
construct only the central action requested by \texttt{Sp1}.
Extend $\rho\circ\mathrm{op}$ complex linearly from the group basis:
\[
 T(z)=\sum_{g\in G}z_g\,\rho(\mathrm{op}([g])).
\]
For $g\in G$ let $C_g=\sum_{h\in G}[hgh^{-1}]$.  If $z$ is central,
its coefficients are constant on conjugacy classes, and counting the terms
of each class gives
\[
 |G|z=\sum_g z_g C_g.\tag{PKT.4}
\]
Each $C_g$ is central and has integer coefficients.  The original matrix
law gives $T(C_gC_h)=T(C_g)T(C_h)$.  Expanding (PKT.4) twice and cancelling
$|G|^2$ proves $T(zz')=T(z)T(z')$ for arbitrary complex central $z,z'$.
The unit, addition and scalar laws follow directly from their original
matrix laws.  Thus $T$ is a complex algebra homomorphism on the center.
If $a\in K[G]$ has coefficientwise image $z$, expansion gives
\[
 T(z)j(x)=j(\mathrm{op}(a)x).
\]
Together with PKT2 this is the full original \texttt{Sp1}/\texttt{Tie1}
pair, including its universal \texttt{cen} clause.  Applying the construction
independently to every already fixed $m_A(b,\mu)$ and $m_L(\mu,j)$ makes
all those ties at once and leaves the matrices themselves unchanged.
% END PKT3

\subsection{The actual finite-type catalogue}

% BEGIN PKT4
\paragraph{PKT4 (finite congruence orbits and the original smooth space).}\label{claim:PKT4}
Let $\mathcal M_\infty$ be the union of the original multiplicity spaces
fixed by positive principal congruence levels, with trivial level character.
It is a vector space: a product of two levels refines both.  Right translation
preserves it, since conjugating a fixed finite adelic matrix and clearing its
finitely many denominators sends a sufficiently deep principal congruence
subgroup into the original one.

Every integral maximal compact subgroup has finite orbits on a vector fixed
at one principal level.  Here is the finite calculation behind this fact.
The integral finite adeles of the totally real field modulo a nonzero
integer have only finitely many classes, since only primes dividing that
integer contribute.  An integral basis of $O_E$ shows that the integral
adelic coefficient algebra for $E$ is a finite module over this ring, so its
quotient is finite as well.  There are only finitely many residue matrices.
Two integral unitary matrices with the same entries modulo a sufficiently
deep level differ, after multiplying by the integral inverse of one, by a
matrix in that principal congruence subgroup.  Their action on the fixed
vector is therefore identical.

The original \texttt{schur}, \texttt{pet}, \texttt{sm}, \texttt{lev} and
\texttt{hc} inputs give a positive Hermitian form on $\mathcal M_\infty$
and finite-dimensional fixed spaces at each principal level.  These are
applications of those explicitly supplied inputs, not new producers for
them.  Their metric is the original $\mathrm{ipM}$.  Right translation is
unitary for this metric: substituting in the same original fundamental-domain
integral preserves its value by right invariance of Haar measure.  In the
source definition by a Schur scalar, the sets of admissible scalars for the
translated and original pair coincide, so the assertion also respects its
fallback convention.  Standard frames and the finite level character imply
that every vector in either of the two original good row spaces lies in
$\mathcal M_\infty$: take a principal subgroup inside the prescribed level
and in the kernel of its finite character.  These arguments preserve the
actual group action and both original levels.
% END PKT4

% BEGIN PKT5
\paragraph{PKT5 (orthogonal projection without a Hilbert completion).}\label{claim:PKT5}
Let a group act unitarily on a positive Hermitian vector space $V$.  Let
$K$ be a subgroup whose orbits on vectors are finite, and suppose
$F=V^K$ is finite dimensional.  For every invariant subspace $U\subset V$
and $x\in F$, there is a unique $u\in U$ such that
\[
 x=u+z,\qquad z\perp U;
 \quad\text{moreover }u,z\in F.\tag{PKT.5}
\]
Project $x$ in the finite-dimensional space $F$ onto $U\cap F$.  Its
residual $z$ is orthogonal to $U\cap F$.  For any $a\in U$, the span
$Y=\langle Ka\rangle$ is finite dimensional and $K$-stable.  The ordinary
orthogonal projection of $z$ onto $Y$ is $K$-fixed by uniqueness and
unitarity.  It lies in $U\cap F$ and hence has zero squared norm.  Thus
$z\perp a$, proving existence.  If two decompositions existed, their
$U$-components would differ by a vector both in $U$ and perpendicular to
$U$, hence would coincide.  This argument requires neither closedness of
$U$ in a completion nor an infinite-dimensional projection theorem.
% END PKT5

% BEGIN PKT6
\paragraph{PKT6 (algebraic irreducible decomposition of every smooth vector).}\label{claim:PKT6}
Under PKT5, every vector of $F$ is a finite sum of vectors in irreducible
invariant subspaces of $V$.  More precisely, prove the assertion inside any
invariant $U$ by induction on $\dim(U\cap F)$.  If this dimension is
positive, choose an invariant $W\subset U$ for which
$0<\dim(W\cap F)$ is least, and put $E=W\cap F$ and
$V_1=\langle GE\rangle$.  A nonzero invariant $W'\subset V_1$ cannot
have $W'\cap F=0$: PKT5 would make $E\perp W'$, and unitarity would then
make $V_1\perp W'$, forcing $W'=0$.  Otherwise minimality gives
$W'\cap F=E$, since it is a subspace of $E$; hence $W'$ contains
$\langle GE\rangle=V_1$.  Thus $V_1$ is irreducible.

For $x\in U\cap F$, PKT5 splits $x$ into its component in $V_1\cap F$
and a component in $U\cap V_1^\perp\cap F$.  The latter fixed space has
strictly smaller dimension because $E\ne0$.  Induction terminates.  Each
vector of the actual $\mathcal M_\infty$ is fixed by one principal level,
so PKT4 makes this argument applicable to it.  The conclusion is algebraic
membership in a sum of irreducibles, not merely density of that sum.
% END PKT6

% BEGIN PKT7
\paragraph{PKT7 (actual function models and orthogonality of distinct types).}\label{claim:PKT7}
An irreducible subspace $W\subset\mathcal M_\infty$ has a model of exactly
the source's type \texttt{FinTyp}.  Choose $w_0\ne0$ and a linear
functional $\ell$ with $\ell(w_0)\ne0$, and send
\[
 w\longmapsto f_w,\qquad f_w(g)=\ell(gw).
\]
The map intertwines the original action with right translation on functions.
Its kernel is invariant and it is nonzero, so irreducibility makes it
injective.  Its image is nonzero, irreducible and smooth, since each $w$ is
fixed by a congruence subgroup.  These are precisely the conditions of the
original finite-type function model.  The inverse on its image embeds that
same model in $\mathcal M_\infty$.

Two such irreducible subspaces of nonisomorphic types are orthogonal.
For each smooth vector of one subspace, PKT5 gives its projection onto the
other.  Uniqueness makes these projections compatible when the fixed level
is refined, linear in the vector, and equivariant under every group element.
A nonzero projection would have zero kernel and full image by irreducibility,
and would therefore be an isomorphism of the two function models.  Hence
the projection is zero.  This proves orthogonality also for the spans of all
copies of two distinct types.  No finiteness of the full adelic group and no
finite-group semisimplicity assertion is used here.
% END PKT7

% BEGIN PKT8
\paragraph{PKT8 (the good cuts preserve the decomposition).}\label{claim:PKT8}
The original good part at one row level is the common kernel of
\[
 T\longmapsto e_0T-T,
 \qquad T\longmapsto r_kT-\chi(k)T\quad(k\text{ in the level}).
 \tag{PKT.6}
\]
Each operator preserves every finite-type isotypic subspace: right
translations do so, and the original $e_0$ is a finite group-algebra sum of
such translations.  If $T$ is good, write it by PKT6 as a finite sum and
combine terms of isomorphic type.  For any one operator in (PKT.6), its
images of the resulting summands still lie in mutually orthogonal type
spaces.  Their sum is zero.  Pairing with each image shows that its squared
norm is zero.  Thus every summand is itself in all the kernels.  This proves
that the actual good part is the sum of its actual isotypic intersections;
it does not require the operators in (PKT.6) to commute.

Each good part is finite dimensional by the original Harish--Chandra input.
Choose a nonzero vector from each occurring distinct type in either row
level.  Their orthogonality bounds their number by the sum of the two finite
dimensions.  The quotient of the union of occurring types by isomorphism
is therefore finite.  Choose one actual occurring representative per class.
The resulting $\Phi$ and $\mathrm{rep}:\Phi\to\mathrm{FinTyp}$ satisfy
all four \texttt{Reps} fields: finiteness, occurrence, uniqueness up to
isomorphism, and exhaustion.  The catalogue is shared by both levels.
% END PKT8

% BEGIN PKT9
\paragraph{PKT9 (the source's conjugate convention and adapted row bases).}\label{claim:PKT9}
The row predicate $\mathrm{isotE}(\pi)$ uses the conjugate finite type;
the lift predicate $\mathrm{isotL}(\pi)$ uses the direct type.  Conjugating
twice gives an isomorphism with the original function model, and conjugating
an intertwiner gives an intertwiner of the conjugate models.  Consequently
PKT8, with representatives chosen for the source's \texttt{OccRows}
convention, gives
\[
 \mathrm{EEc}_b=\sum_{\pi\in\Phi}\mathrm{prtE}_b(\mathrm{rep}\,\pi),
 \qquad
 \mathrm{prtE}_b(\pi)\perp\mathrm{prtE}_b(\pi')\quad(\pi\not\cong\pi').
 \tag{PKT.7}
\]
The tensor-lift identification of multiplicity maps is a linear equivalence,
so these equalities and the orthogonality hold in the actual ambient spaces
used by \texttt{ChT}.  Similarly the actual lift parts
$\mathrm{prtW}_i(\pi)$ have zero intersection for distinct catalogue indices.
If $j_E:E\to\mathrm{EEc}_b$ is an isometric isomorphism onto that subspace,
pull back (PKT.7), choose an orthonormal basis in each part and concatenate
the finitely many bases.  Label a basis vector by its part.  The span of the
vectors with label $\pi$ is exactly $j_E^{-1}\mathrm{prtE}_b(\mathrm{rep}\,\pi)$.
Empty parts contribute no vectors.  This proves the full \texttt{ft} and
\texttt{disj} clauses, including zero-dimensional row spaces.
% END PKT9

\subsection{Definite coefficients, class sums and both chains}

% BEGIN PKT10
\paragraph{PKT10 (the actual definite compact models and their dual pairing).}\label{claim:PKT10}
At the terminal definite row stage the original sign tuple is all positive;
at a definite lift coefficient stage it is uniformly of the prescribed sign.
The original lowest compact type is nonzero and finite dimensional.  Since
in these definite cases the maximal compact real group is the entire real
group, that lowest type is invariant under the entire group.  It is closed,
and irreducibility of the original unitary model makes it the whole model.
In particular the smooth vectors equal the whole finite-dimensional space.
This also covers rank zero, where the original character model has dimension
one.  These conclusions use the irreducibility and lowest-type assertions
already included in the original representation data.

The tensor product over the real places is irreducible for their product
group.  One elementary justification is useful here.  The algebra spanned
by a finite-dimensional irreducible unitary representation is a unital
$*$-algebra.  Spectral projections of its selfadjoint elements belong to the
algebra, being polynomials in those elements.  Choose a nonzero projection
of least rank.  Its corner algebra consists of scalars, since otherwise a
selfadjoint corner element gives a smaller spectral projection.  If its rank
were at least two, two orthogonal vectors $u,v$ in its range would satisfy
$\langle Au,Bv\rangle=0$ for all algebra elements $A,B$.  Both cyclic spans
are the entire irreducible space, a contradiction.  Thus it contains a
rank-one projection; multiplying it on both sides spans all rank-one
operators, so the algebra is the full matrix algebra.  The product group
therefore spans the tensor product of the full matrix algebras, which is the
full endomorphism algebra of the tensor product.

The original pairing $\mathrm{prX}$ between the coefficient space and this
tensor model is equivariant and nonzero.  Its kernels are invariant, hence
zero by the two irreducibilities.  Since both spaces are finite dimensional,
the two injections into their duals give equality of their dimensions and
a perfect pairing.  Its actual map is consequently an isomorphism onto the
dual tensor space; it is not replaced by an unspecified pairing.  In
particular, with $d_X=\prod_\tau\dim M_{X,\tau}$,
\[
 \dim W_X=d_X>0.\tag{PKT.8}
\]
The original \texttt{PrIso} normalization says that the Riesz vectors $a_v$
of pure tensors obey
$\langle a_{v'},a_v\rangle=\prod_\tau\langle v_\tau,v'_\tau\rangle$.
This follows directly by expanding in orthonormal bases; the reversed order
on the right records the conjugate linearity of the Riesz identification.
% END PKT10

% BEGIN PKT11
\paragraph{PKT11 (the original compact forms give all multiplicity maps).}\label{claim:PKT11}
For a finite adelic coefficient form $F$ in the original definite form space,
the original extension and pairing give
\[
 T_F(v)(g)=\mathrm{prX}\bigl(v,\rho(g^{-1})F(g_f)\bigr).
 \tag{PKT.9}
\]
The prescribed rational equivariance of $F$ proves left rational invariance;
its finite level law proves the same finite character law for $T_F$.
The compact real action and equivariance of $\mathrm{prX}$ prove the exact
real intertwining law.

For completeness these functions have all the required automorphic
regularity, not just the transformation law.  A continuous finite-dimensional
representation of a real Lie group is smooth: choose a smooth approximate
identity supported near $1$.  Its averaged operator tends to the identity
and is therefore invertible for a sufficiently small support; its range
consists of smooth vectors by differentiating the smooth convolution kernel.
Thus every vector is smooth.  The derivatives of (PKT.9) are matrix
coefficients with the corresponding fixed products of Lie matrices.
The finite class set and unitarity give a bound for each such derivative,
uniform in $g$.  Hence the uniform moderate-growth condition holds with
exponent zero.  The center acts on the finite-dimensional coefficient
function space, so its orbit spans a finite-dimensional space.  A definite
rational Hermitian space has no nonzero isotropic rational flag, so the
cuspidal conditions along proper rational isotropic parabolics are vacuous.
Finally the multilinear maps in the finite-dimensional tensor variables are
continuous; a fixed coordinate expression bounds them by a constant times
the product of the norms.  This supplies the source's Sobolev-continuity
condition using the order-zero seminorm.  Thus (PKT.9) is the actual
\texttt{IsMultMap}, including its regularity clauses.

Conversely let $T$ be in the original definite multiplicity space.  At each
$g$, the functional $v\mapsto T(v)(g)$ has a unique coefficient vector
$\Psi(g)$ by PKT10.  Real equivariance gives
$\Psi(ga)=\rho(a^{-1})\Psi(g)$, first for one real factor and then for their
product.  Therefore $\rho(g)\Psi(g)$ depends only on $g_f$.  Call it $F(g_f)$.
Left rational invariance and the finite character law of $T$ are exactly the
rational and level laws of this $F$.  The same formula gives $T=T_F$.
The two constructions are inverse.  This proves the original
\texttt{TopMultH} on each prescribed tuple, with the original coefficient
extension, pairing and level subgroup.
% END PKT11

% BEGIN PKT12
\paragraph{PKT12 (Schur unfolding with the literal normalization).}\label{claim:PKT12}
Let $\Omega$ be the preimage of the prescribed finite level in the definite
adelic group.  The original measure normalization gives $\nu_X(\Omega)=1$.
The representation is trivial on finite-only elements, while $\Omega$
contains the whole compact real group.  Thus its action on $W_X$ is unitary
and irreducible.  Averaging a rank-one operator over $\Omega$ gives an
operator commuting with this action.  Irreducibility makes it a scalar;
its trace determines that scalar.  Equivalently,
\[
 d_X\!\int_\Omega
 \overline{\langle a,\rho(k^{-1})u\rangle}
 \langle a',\rho(k^{-1})u'\rangle\,d\nu_X(k)
 =\langle a',a\rangle\langle u,u'\rangle .
 \tag{PKT.10}
\]
All integrands here are continuous on the compact quotient acting on the
finite-dimensional coefficient space, and bounded; their integrability
requires no cuspidal estimate.

Take the actual finite class representatives $y_c$ from the original class
presentation and their finite rational stabilizers $\Gamma_c$.  The value
at $y_ck$ is $\chi(k)\rho(k^{-1})F(y_c)$.  In a product of one conjugate
and one unconjugated coefficient, the unit-modulus character cancels.
Unfolding the original fundamental domain yields the sum of the integrals
on these finitely many pieces with weights $1/|\Gamma_c|$.  To see the
weight directly, the map from a level coset to the rational quotient has
exactly $|\Gamma_c|$ preimages on that piece; partition a fundamental
domain for this finite action and apply translation invariance to its
translates.  No measure is renormalized during this step.  Applying
(PKT.10) to each piece gives the original weighted class-sum pairing.
Using the Riesz identity in PKT10, and a nonzero pure tensor to determine
the Schur scalar uniquely, gives
\[
 d_X\,\mathrm{ipMT}\bigl(T_h,T_{h'}\bigr)
       =\langle h,h'\rangle_{H_X}.
 \tag{PKT.11}
\]
In particular the scalar exists, so the fallback branch of the definition
of $\mathrm{ipMT}$ is not used in this computation.
% END PKT12

% BEGIN PKT13
\paragraph{PKT13 (exact level spans and both coordinate Gram identities).}\label{claim:PKT13}
Use the original rational class-value realization MR1 and its character
averaging in MR2.  For a finite subgroup $P$ and character $\psi$, the
operator
\[
 \Pi_\psi=|P|^{-1}\sum_{p\in P}\psi(p)^{-1}r_p
 \tag{PKT.12}
\]
is the identity on the desired character space and maps the whole form
space onto it.  This is checked by changing the summation variable under
$r_q$; applying it twice then gives the same operator.  The coefficient
field contains the actual character values.  The original action tie says
that the rational coordinate action and adelic translation in (PKT.12)
agree.  Applying (PKT.12) to a rational spanning family therefore proves
that the actual principal-level adelic forms are spanned by the original
\texttt{rowForm} or \texttt{liftForm} images.  This is equality of ranges,
not just equality of dimensions.

The original normalized coordinate Gram matrix is
$|P|^{-1}E^*HE$, where $E$ is the actual cut-basis matrix and $H$ is the
weighted class-sum Gram matrix.  Thus MR1 and TL7 give, for original row
coordinates $x,z$,
\[
 \langle j_X\mathrm{rowForm}(x),j_X\mathrm{rowForm}(z)\rangle
 =n_b\,\sigma\bigl(\mathrm{hform}(m_A.H)(x,z)\bigr).
 \tag{PKT.13}
\]
Combining this with (PKT.11) and the literal definition
$\mathrm{ipXP}_b=(d_X/n_b)\mathrm{ipMT}$ proves that the coordinate-to-top
map preserves exactly $\sigma(m_A.H)$.  The factor is $n_b$, with $b$
unchanged throughout.

For lift-source coordinates the original coefficient extension
$\mathrm{extV}$ is injective and the source inner product is its transported
coefficient inner product divided by $n_0$.  The same argument at
$(P_0,\psi_0)$ proves
\[
 \mathrm{ipA}_{\mu,i}
  (\mathrm{extV}\,j_V\mathrm{liftForm}(x),
   \mathrm{extV}\,j_V\mathrm{liftForm}(z))
 =\sigma\bigl(\mathrm{hform}(m_L.H)(x,z)\bigr).
 \tag{PKT.14}
\]
The range is exactly the original intersection called $\mathrm{Amt}$:
intersect the coefficient-extension range with the finite level law, then
apply (PKT.12) before the injective extension.  These statements retain the
original class weights, the original $\psi_0^{-1}$ target convention, and
both distinct normalization factors.
% END PKT13

% BEGIN PKT14
\paragraph{PKT14 (Hilbert models and the three Riesz transitions).}\label{claim:PKT14}
The original inputs make $\mathrm{ipE}_b$, $\mathrm{ipP}_{b,\mu,j}$ and
$\mathrm{ipXP}_{b,\mu}$ positive on their stated original multiplicity
spaces.  Indeed the original finite-level Schur--Petersson assertion applies
to the same Haar measure and fundamental domain; division by $n_b>0$
preserves positivity.  At the top, the additional factor $d_X$ is positive
by PKT10.  The original base and intermediate subspaces satisfy
\[
 \mathrm{EEc}_b\subset\mathrm{CC}_{b,\mu,0},\qquad
 \mathrm{CCP}_{b,\mu,j}\subset\mathrm{CC}_{b,\mu,j},\qquad
 \mathrm{CCXP}_{b,\mu}\subset\mathrm{CCX}_{b,\mu}.
 \tag{PKT.15}
\]
The first follows because the base row rank is $r$ and its lower block has
rank zero; the others express the stronger character-level condition inside
the corresponding full level.  They are the actual subgroup inclusions.

Make $E_b$ the finite-dimensional Hilbert model of $\mathrm{EEc}_b$, with
isometric inclusion $j_E$.  For each used tuple and $1\le j\le m-1$ make
$M_{\mu,j}$ the analogous model of $\mathrm{CCP}_{b,\mu,j}$, with $j_M$.
At $j=0$ use $E_b$, and at $j=m$ use the already fixed top $s_A(b,\mu)$.
PKT13 and the complex basis of PKT2 define its unique isometric map $j_X$
by the original row-coordinate formula.  Its range is $\mathrm{CCXP}$.
The inverse $y.\mathrm{tupT}$ is used only on tuples in the image of
$y.\mathrm{tup}$; every used tuple is in that image, including $\mu_0$
by $y.\mathrm{tup0}$.

For a finite-dimensional Hilbert model $H$ embedded by $j$ in an ambient
pairing space, the functional $x\mapsto \mathrm{ip}(jx,w)$ has a unique
Riesz vector $r(w)$ with $\langle x,r(w)\rangle=\mathrm{ip}(jx,w)$.
Expansion in an orthonormal basis shows that $r$ is linear in $w$.
Define the successive maps by applying this construction to the actual
$\mathrm{ioC}$, and at the last step to the actual $\mathrm{ioX}$.
At $j=0$ use the $E_b$ pairing, at $1\le j\le m-2$ the corresponding
intermediate pairing, and at $j=m-1$ the penultimate pairing with $j_X$
on the input.  This gives exactly the three displayed adjoint-pairing
identities of \texttt{ChT}.  The equality $m-1+1=m$ accounts for its
terminal heterogeneous-equality cast.  Unused tuples receive zero transition
maps.  Outside the intermediate range the arbitrary auxiliary models can
be $E_b$; they do not change any specified endpoint or transition identity.
% END PKT14

% BEGIN PKT15
\paragraph{PKT15 (both complete original chain comparisons).}\label{claim:PKT15}
Perform PKT14 for $b=\mathrm{false}$ and $b=\mathrm{true}$ with the shared
catalogue PKT8.  In each base Hilbert space choose the adapted orthonormal
basis from PKT9.  Define \texttt{multZero} and \texttt{multTop} by the
identity identifications in the piecewise definition of $M_{\mu,j}$.
The resulting \texttt{ChD} and \texttt{ChT} have all their fields:
$j_E$ has the original range and metric; the $j_M$ and $j_X$ have the
original ranges, injectivity and metrics on every used tuple; all three
$\mathrm{io}$ identities are PKT14; unused transitions vanish; and the
labelled basis and distinct-part intersections are PKT9.  Finally the
coordinate definition of $j_X$ gives, for every original $t$ and $x$,
\[
 j_X(y.\mathrm{tup}(t))(j_A(t)x)
 =R.c_X.\mathrm{toMap}(R.\mathrm{prX})
       \bigl((R.H_X(t)).e(R.j_X(t)(\mathrm{rowForm}(x)))\bigr).
 \tag{PKT.16}
\]
Equality on the fixed complexified coordinate basis extends this formula
to all $x$.  This is the complete original \texttt{top} field, on the
already chosen $s_A$, not a later independently chosen top model.
% END PKT15

\subsection{The lift lattice and its conditional theta map}

% BEGIN PKT16
\paragraph{PKT16 (the exact good integral carrier in the target model).}\label{claim:PKT16}
Fix a target index $i$.  The original \texttt{Holds.aux.integral} supplies
the $K_s$-rational form, the full finite $O_{K_s}$-lattice
$M_K=(R.FW_i).MK$, and its stability under the original finite group.  The retained coefficient-field
and character conditions imply that the actual good cut
\[
 W_{\rm good}=W_{\rm good}(P_0,\psi_0^{-1},e)
\]
is defined over $K_s$.  The integral-descent argument DI17, with these
literal inputs, gives
\[
 O_K\text{-span}(M_K\cap W_{\rm good})
   =(O_K\text{-span}M_K)\cap W_{\rm good},
 \tag{PKT.17}
\]
with the field extensions understood exactly as in the original
\texttt{FW} model.  The relevant intersection is a finite $O_{K_s}$-module because
it is a submodule of the finite $O_{K_s}$-lattice over a Noetherian ring.  It is full
in the good subspace: choose a rational basis of that subspace and clear
the finitely many denominators in the ambient full lattice.

Suppose $E:M\to H_W$ is injective and has range
$W^{\rm at}_i.\mathrm{comap}(H_W.e)$.  The original level equality in
\texttt{Holds.aux.levelW} identifies this subspace with the complex span
of $(R.FW_i).j(\mathrm{sg0})(M_K\cap W_{\rm good})$.  Choose integer
generators of the latter integral intersection and lift their images
uniquely through $E$.  The $O$-lattice they generate in $M$ then satisfies
\[
 x\in\Lambda_i\ \Longleftrightarrow\
 \exists v\in M_K\cap W_{\rm good},\quad
 E x=(R.FW_i).j(\mathrm{sg0})v.\tag{PKT.18}
\]
Both implications follow from the chosen integer generators and the
integrality of the original $O$-action, as in PKT2.  Specifically, its
image lies in $K_s$ by PKT1; integrality in $K$ descends along the
field embedding to integrality in $K_s$, so it acts through $O_{K_s}$.
Full complex spanning
follows from (PKT.17) and the exact range of $E$.  Thus no arbitrary
commensurable lattice is used.
% END PKT16

% BEGIN PKT17
\paragraph{PKT17 (the target Hilbert space with its factor $n_0$).}\label{claim:PKT17}
For $1\le j\le m$ put $i=j-1$ and
\[
 S_i=(R.W^{\rm at}_i(\mathrm{sg0},\mathrm{levL}_0))
                   .\mathrm{comap}(R.HW_i(\mathrm{sg0}).e).
\]
It is finite dimensional: $HW_i$ is the original finite good coefficient
model and its embedding into multiplicity maps is injective.  Choose an
orthonormal basis of $S_i$, set $M_j=\mathbb C^{\dim S_i}$, and let
$E_i:M_j\to HW_i$ be $\sqrt{n_0}$ times the basis inclusion.  Then
\[
 \operatorname{range}E_i=S_i,\qquad
 \|E_i x\|^2=n_0\|x\|^2.\tag{PKT.19}
\]
Define $j_{M_j}$ to be the actual tensor lift of $HW_i.e\circ E_i$.
The original $HW$ metric identity and the definition of $\mathrm{ipWc}$
give
\[
 \mathrm{ipWc}_i(\operatorname{lift}(HW_i.eu),
                   \operatorname{lift}(HW_i.ev))
       =n_0^{-1}\langle u,v\rangle.\tag{PKT.20}
\]
To check the constants, its defining $\mathrm{ipAt}$ contributes the
product of the local lowest-type dimensions, while the original $HW$
normalization divides by that same product; scaling the original measure
by $n_0^{-1}$ contributes the remaining factor.  No other scalar survives.
Thus $j_{M_j}$ is an isometry onto exactly $\mathrm{Wc}_i$.  Equip $M_j$
with the lattice from PKT16.  At $j=0$ and $j>m$ use the zero Hilbert
space and its zero lattice.  These choices give all target norm, range,
lattice, injectivity and out-of-range clauses in the original \texttt{LfT}.
% END PKT17

% BEGIN PKT18
\paragraph{PKT18 (the prescribed source embeddings and the conditional lift).}\label{claim:PKT18}
For a used tuple $\mu=y.\mathrm{tup}(t)$, send the complexified coordinate
basis of $s_L(\mu,i+1)$ to
\[
 \mathrm{extV}_{i,t}\,j_{V,i,t}\,\mathrm{liftForm}(e_k).
\]
PKT13 says that this extends to an injective isometry $j_{M_t}$ with
range exactly $\mathrm{Amt}_{\mu,i}$.  The inverse tuple $t$ is unique
because $y.\mathrm{tup}$ has the given left inverse.  The map is set to
zero outside the image when it has no required source comparison.

Now let $U_i$ be the subspace of original $\mathrm{MultW}_i$ consisting
of maps $w$ with $e_W\circ w\in\mathrm{Wc}_i$.  Restriction
\[
 r:U_i\longrightarrow\operatorname{Hom}_{\mathbb C}
       (\text{lowest compact type},\text{functions}),
 \qquad r(w)=w\circ\mathrm{embK}_i
 \tag{PKT.21}
\]
is injective.  Indeed each original local lowest type is nonzero; choose a
nonzero vector in each and apply the one-lowest-tuple determination argument
DI4 to the difference of two multiplicity maps.  All its original
intertwining and smoothness conditions are already part of membership in
$\mathrm{MultW}_i$.

Define a linear left inverse to $r$ on its image and extend it to the full
ambient vector space by a basis extension.  Compose this map with the
original $\mathrm{thC}_{\mu,i}\circ j_{M_t}$, then with
$w\mapsto e_W\circ w$, and finally with the inverse of $j_{M_j}$ on
$\mathrm{Wc}_i$.  Call the resulting linear map $\theta_{\mu,i+1}$.
If an original $w\in\mathrm{MultW}_i$ satisfies
\[
 w\circ\mathrm{embK}_i=\mathrm{thC}_{\mu,i}(j_{M_t}x),
 \qquad e_W\circ w\in\mathrm{Wc}_i,
\]
then the left inverse recovers that very $w$, and
\[
 j_{M_j}(\theta_{\mu,i+1}x)=e_W\circ w.\tag{PKT.22}
\]
This proves precisely the conditional theta clause.  It does not assert
that such a $w$ exists for every $x$, and it does not consume or replace
the independent theta-extension producer.  Set all $\theta$ maps on unused
tuples and outside the valid index range to zero.
% END PKT18

% BEGIN PKT19
\paragraph{PKT19 (the complete original lift comparison).}\label{claim:PKT19}
Set the target type subspace at $j=i+1$ to
$j_{M_j}^{-1}\mathrm{prtW}_i(\mathrm{rep}\,\pi)$.  PKT9 gives the
required distinct-part intersections.  Use the zero subspace outside the
valid range.  PKT16--PKT18 now give every \texttt{LfD}/\texttt{LfT}
field on the fixed $s_L$: the original $e_W$ with its range, factor $n_0$
and exact integral carrier; the tensor-lift reading and its isometry onto
$\mathrm{Wc}$; the source $j_{M_t}$ with its exact range and metric; the
conditional theta equation; unused-tuple vanishing; out-of-range zero
spaces; and all type predicates.  The coordinate construction also gives
for every original $t$ and $x$
\[
 j_{M_t}(y.\mathrm{tup}(t),i)(j_L(i,t)x)
       =R.\mathrm{extV}(i,t)
          (R.j_V(i,t)(\mathrm{liftForm}(A_3.A,P_0,\psi_0,i,t,x))).
 \tag{PKT.23}
\]
This is the literal final \texttt{read} clause, using the same $j_L$
from PKT3 and (PKT.1).
% END PKT19

% BEGIN PKT20
\paragraph{PKT20 (sharp row coordinates and the exact loss parameter).}\label{claim:PKT20}
Apply the independently proved sharp coordinate theorem RDR7 to the actual
$j_E$ of the true-level chain in PKT15.  It gives $n_R$, a complex basis
$c_R$ of that same $E_{\rm true}$, and every original all-place frame
$X_o(v)$, with
\[
\begin{split}
 j_E\!\left(\sum_a\sigma(x_a)c_R(a)\right)
 &=\operatorname{lift}\bigl(R.HE.e\,
      \mathrm{Tw.out}(i_C(\mathrm{cpt}(x)))\bigr),\\
 X_o(v)&=X_E(v)\text{ reindexed by the same dimension equality},
 \qquad \det X_o(v)\ne0.
\end{split}\tag{PKT.24}
\]
Here $i_C$, the involutively semilinear $\mathrm{cpt}$, the distinguished
embedding, and $X_E$ are exactly those of $(R.\mathrm{toCx3}(F))$.
This is the full original \texttt{RrT}, including all embeddings, and not
only a dimension comparison.

Let $D_{\rm all}$ be the dimension of the full original row multiplicity
space at $(P_1,\psi_1)$, before the good idempotent cut, and let
$n=\dim E_{\rm true}$.  If $D_{\rm all}=0$, the full multiplicity space
is zero; its good subspace and hence $E_{\rm true}$ are zero.  Define
\[
 \delta_p=\begin{cases}0,&D_{\rm all}=0,\\
                  1-n/D_{\rm all},&D_{\rm all}>0.
          \end{cases}\tag{PKT.25}
\]
Then $\delta_p\le1$ and $n=(1-\delta_p)D_{\rm all}$ in both cases.
Only this upper bound on $\delta_p$ is required and proved here; no
additional nonnegativity assertion is included in the packet.
% END PKT20

% BEGIN PKT21
\paragraph{PKT21 (full original \texttt{exists\_pkT}).}\label{claim:PKT21}
Under exactly PKT1's original premises there is $p:\mathrm{PkD}(D,G)$
with $p.m_L=m_L$ and a member of $\mathrm{PkT}(F,y,m_A,p)$.
Take the finite catalogue PKT8, every complexified space and tie PKT3,
both chains PKT15, the lift PKT19, and the numerical and row-coordinate
data PKT20.  The laws and equalities (PKT.1) are the given ones, copied
without alteration.  The $\mathrm{PkT}$ fields \texttt{rep/reps},
\texttt{tA}, \texttt{lawsA/lawsL}, \texttt{hA/hL}, \texttt{ch},
\texttt{lf}, \texttt{dall}, and \texttt{rr} are respectively the exact
certificates just proved.  All universal choices are made for the same
fixed family before constructing $p$; the dependent lift and chain
coordinates use the resulting same choices.  This proves the original
endpoint without any new representation-catalogue, lattice, pairing or
intertwiner premise.
% END PKT21

\subsection{Construction of the entire original row package}

% BEGIN PKT22
\paragraph{PKT22 (all original $\mathrm{Hd}$ fields).}\label{claim:PKT22}
Now fix the original $B,f,T,w,Z_0,h_X,W_t$ and a given realization
$z_0:\mathrm{Dat}(w.\mathrm{st}(Z_0),w.I,h_X,T,W_t)$.
Assume the original host equality $F=E^+$, the standard self-dual lattice
conditions for $Y_0$ and every $W_i$, and the original inputs and holds
conditions.  We construct the additional \texttt{Hd} data rather than
assuming compatible frames.

At each finite place the original unramified CM condition and the given
self-dual lattice give an integral orthonormal frame by DR1, including at
residue characteristic two.  Relative to the prescribed rational lattice
frame, both the change matrix and its inverse have integral entries.
These local matrices glue to a finite adelic matrix: a fixed integral basis
of $O_E$ expresses each local integral entry as a tuple of integral
coefficients over the totally real completion; those coefficient tuples
are integral finite adeles.  Glue the inverse entries in the same way.
The local identities $AB=BA=1$ and $A^*HA=1$ imply the adelic identities
because evaluation at all finite places is injective.  This constructs a
finite adelic orthonormal frame, not a claimed free global $O_E$-basis of
the lattice.

Conjugation by that frame identifies the standard finite unitary group
with the original finite-only group, and identifies their maximal compact
subgroups.  Let these homomorphisms be $P_Y$ and $P_{W_i}$, and set
$\iota_{W_i}=P_{W_i}P_Y^{-1}$.  The inverse frame matrices show directly
that this is the original \texttt{IsIsomFin} with its conjugation and
lattice clauses, and that \texttt{FrameMatch} holds with parameter one.
Thus the frame maps and their standardness are simultaneous original maps.

The used-tuple inclusion $\mathrm{ixOf}$ is injective.  Define its left
inverse by its unique preimage on the image, and by the base tuple elsewhere.
It sends the original base tuple correctly.  The rank assumption gives
$m\ge2$, hence $m>0$.  PKT11 gives \texttt{TopMultH}; injectivity of
$\mathrm{extV}$ follows from injectivity of the original coefficient
extension and of the given coefficient model $HV.e$.  These facts, with
the constructed frames, are every field of \texttt{Hd}, with
$y.\mathrm{tup}=\mathrm{ixOf}$.
% END PKT22

% BEGIN PKT23
\paragraph{PKT23 (all original coordinate laws and restriction identities).}\label{claim:PKT23}
Use the original all-tuples matrices $m_{AU}(b,a)$ and $m_{LU}(a,j)$.
The latter is the actual lift matrix at $i=j-1$ for $1\le j\le m$,
and is \texttt{Mat1.zero} otherwise.  We verify their laws from their
original finite coordinates.

Let $E_P$ be the chosen coordinate inclusion of the character cut, and
$L_P$ its left inverse.  The cut's Gram matrix is
$|P|^{-1}E_P^*HE_P$ and its integral lattice is the inverse image of the
original coefficient lattice under $E_P$.  Positivity follows from TL7
and injectivity of $E_P$.  Fullness follows by clearing denominators of
a basis in the full ambient lattice; finite generation follows by
Noetherianity.  This proves the literal lattice law, not just discreteness
at the distinguished embedding.

For a group-algebra element $a$, the matrix action is
$L_P\rho(a)E_P$.  It is additive, scalar-linear, and sends $1$ to $1$.
If $a$ is central then $\rho(a)$ commutes with the character projector
(PKT.12), so it preserves the cut.  Consequently
$E_PL_P\rho(a)E_P=\rho(a)E_P$, which proves multiplicativity for two
central elements.  The original good-character catalogue in $z_0.\mathrm{good}$
consists of absolute irreducible finite-group characters.  Averaging the
usual character matrix coefficients gives orthogonal central idempotents;
the sum over the specified catalogue is therefore the original central
idempotent.  The finite action is unitary for the actual weighted class-sum
form TL5--TL7, so this idempotent is selfadjoint.  Its compressed matrix is
exactly the stored $E$.  These are all ten \texttt{Mat1.Laws} fields.
For the zero matrix record the space has dimension zero; every identity
and the full-lattice assertion are vacuous.

The restricted family at a used tuple $t$ and the all-tuples family at
$\mathrm{ixOf}(t)$ have identical character, weight, level and class-value
data.  Their coordinate definitions therefore use the same cut basis,
Gram matrix, lattice and action.  Substituting these data gives the exact
identities (PKT.1), rather than merely isomorphic coordinate records.
Finally $z_0.K_s\mathrm{\_le}$ supplies the coefficient-ring hypothesis
of PKT1.  Thus PKT21 applies to these actual all-tuples matrices.
% END PKT23

% BEGIN PKT24
\paragraph{PKT24 (actual period coordinates and the full $\mathrm{PerT}$).}\label{claim:PKT24}
For every tuple $a$, the original complex row-form space
$\mathrm{MtCx}(w.I,a)$ has the complexified rational basis of MR1--MR2.
Map that basis to the same-coordinate basis of $s_A(\mathrm{true},a)$
from PKT3.  This gives a linear equivalence $e_a$ with the exact
\texttt{IsCrd} identity for the original \texttt{rowForm}.

For every word $u$ and tuple $a$, there is an explicit equivalence
$g_{u,a}$ from this row-form space to the original period-form space.
Here is its descent argument.  Pull a row function back along the original
surjection from framed points to unframed row points.  Rational and finite
level equivariance become exactly the period frame law.  Conversely two
framed points with the same unframed point differ by an element in the
kernel of frame reduction.  At the disjoint sets $vG_s$ and $Q(a)$ the
local changes can be chosen independently.  Their assembled reduction is
the identity, their rational part is the identity, and the product of their
level characters is one.  The frame law therefore makes the two function
values equal.  The function descends uniquely.  Equivariance of the
surjection proves the descended row law.  The disjointness used here is
exactly the original selection condition in $w.\mathrm{ok}$; the
reduction maps and inverse characters are the ones in
\texttt{RightTranslationActionApply}.  The two pullback/descent operations
are inverse and complex linear.

Set $d.e_0(u)=e_{\mu_0}g_{u,\mu_0}^{-1}$ and
$d.e_T(u,a)=e_ag_{u,a}^{-1}$.  These maps give the original period record
$\mathrm{PerD}(p)$.  Their inverses undo the same $g_{u,a}$, so the
\texttt{SameForm} equalities hold pointwise at every original frame, for
both the reference and twisted periods.  Together with the initial
coordinate identity they give all \texttt{PerT} fields \texttt{e},
\texttt{ref}, \texttt{tw}, and \texttt{crd}.  No statement about the value
or nonvanishing of a period integral was needed.
% END PKT24

% BEGIN PKT25
\paragraph{PKT25 (the prescribed auxiliary rows, weights and Sha dictionary).}\label{claim:PKT25}
The auxiliary $W_0:\mathrm{RowsRest}(B,w.C)$ can be chosen independently
of the comparison packet while retaining the original column-period spaces.
Indeed MR1, the original integral form lattice and its positive metric give
for each tuple an isometric Euclidean model and an $O$-lattice with exactly
the original integral carrier; the construction is PKT2 applied to an
integral basis and its Gram matrix.  Transport the equivalences $g_{u,a}$
of PKT24 to these models.  Put $\mathrm{Mu}_u$ equal to the original
ramified column space, with its original positive inner product and the
identity identification $e_\mu$.

For clarity the permitted auxiliary fields are explicit.  Take $EE=EE_0=0$,
$D_{\rm rows}=D_0=D_{\rm all}=\delta_p=0$, the empty type set and empty
orthonormal row bases.  Let $\mathrm{Amb}_a$ be the Euclidean model just
chosen; let $\mathrm{Mult}_{a,j}=0$ except at $j=m$, where it is
$\mathrm{Amb}_a$.  All intermediate maps are zero, the ambient lattice
is the chosen one, and its idempotent is the identity.  Define the column
period by the actual $P\widetilde{\mathrm{Col}}$, using the inverse
reference equivalence at $\mu_0$ and the inverse twisted equivalence and
original $e_{0,\rm Tw}$ elsewhere.  Initially take $B_u=0$ and $z_e=1$.
Take $M_j=M_{t,\mu,j}=0$, their zero lattices, identity idempotents,
zero theta maps, and zero $\lambda$, $I_j$, $j_w$ and $\eta_j$; use the
full subspace as each target type subspace.  Every zero-space basis and
idempotence clause then holds directly.  The reference and twisted period
ties follow by the two branches in the definition, and $e_\mu$ preserves
the original ramified form.  These are initial stand-in fields only: the
later operation $p.\mathrm{what}_2$ inserts the actual chains and lifts
of PKT21.

The weight fields are not arbitrary.  Define $s_i$ to be the original
$v_i$ for $i<k-1$, then the original $-b_{i-(k-1)}$ when in range, and
zero otherwise; define $e_j=e_{j-1}^{\rm orig}$ for $1\le j\le m$ and
zero otherwise.  Put
\[
\begin{split}
 a^+_{j,i}&=s_i+e_j-(n_V-1-2i)/2-(k+j-1-m+j)/2,\\
 a^-_{j,c}&=(b_{j+c}\text{ if in range, else }0)-e_j-(t+1+2c)/2.
\end{split}\tag{PKT.26}
\]
All original $v,b,e$ are integers.  Writing $t=2s+1$ and using the
original rank identity reduces the two expressions to
$a+b-(k+j-2-i)-(s+1)$ and $a-b-(s+1+c)$ respectively, with integer
$a,b$.  Thus set $\mathrm{kaBig}=a^+$ and
$\mathrm{kaSmall}=-a^-$.  These equalities prove the exact
\texttt{KapDef} for every later $V,S,N,G$, since these stages do not
alter the weight formulas.

Choose $\mu_{\rm Tw}(v)$ as a function of the character
$T.\xi_{\rm Tw}(f,v)$ in the original image of $w.I.\xi$; equal
characters receive equal preimages.  Conjugate places have equal characters,
so the tuple is conjugation-invariant.  Set
$\mu_{\rm At}(v)=\mu_{\rm Tw}(v)$ on $W_t$ and $\mu_0$ off $W_t$;
the assumed conjugation-invariance of $W_t$ gives all \texttt{MuW} clauses.
Finally take the actual $\mathrm{shaTrue}$ from LS18: its value is the
natural conversion of the original extended module length.  When that
length is a finite integer $n$, its conversion is $n$, exactly
\texttt{ShaOK}; no finiteness of all those lengths is asserted.  This
completes all prescribed auxiliary fields.
% END PKT25

% BEGIN PKT26
\paragraph{PKT26 (one numerical record valid for every later stage).}\label{claim:PKT26}
Fix the actual $p,d,W_0$ just constructed.  Temporarily assemble them with
zero numerical fields and zero $\eta$, and call the resulting underlying
row data $D^0$.  Its spaces, periods, sections, projections, theta maps and
integral generators are independent of these numerical placeholders.  Set
\[
\begin{split}
 B_u(\phi,x)&=\mathrm{perCol}_{u,\mu_0}(\phi,
                                      \mathrm{sec}_{\mu_0}x),\\
 \lambda_{\mu,j,\pi}
 &=\frac{\Re\operatorname{tr}
   (\theta_P\theta_T\operatorname{pr}_{\mathrm{isot}(j,\pi)})}
          {\dim\mathrm{isot}(j,\pi)}.
\end{split}\tag{PKT.27}
\]
Use the source's division-by-zero convention for a zero type space.  The
first expression is bilinear since the actual period and section are
linear.  These definitions are precisely the \texttt{bu} and
\texttt{lam} clauses of \texttt{NumOK}.

A scalar satisfying the source's period factorization, together with a
nonzero period witnessing that factorization, is unique: evaluating two
factorizations at that witness and cancelling the nonzero remaining factor
gives equality.  At $\mu_0$ the identity
$\mathrm{sec}\,\mathrm{iota}=1$ on the projected subspace makes it one.
Define $z_e(u,\mu_0)=1$; at another tuple choose a unit of the original
field giving such a factorization if one exists, and choose $1$ otherwise.
If the \texttt{NumOK.ze} antecedent supplies a factorization and a field
value, its nonzero witness proves that field value is a unit, and uniqueness
identifies it with this choice.  Thus the required conditional clause holds
without asserting existence of any new factorization or nonzero period.

For the two index fields let $T_{j,\mu}$ be the $O_K$-span of the original
rational vectors whose images belong to the original $T_j$ set.  Its dual
consists of vectors in its $K$-span pairing integrally with all of
$T_{j,\mu}$.  Define $j_w$ as the original local index of these two
modules.  In the complex target let $L_j^{\rm can}$ be the $O_K$-span of
the actual target-lattice generators, and let $T_j^{\rm can}$ consist of
the theta images of $T_{j,\mu}$ which belong to the $K$-span of those
same generators.  Define $I_j$ as their local index.

These indices agree with those computed in every subsequently supplied
rational target model $Y$.  Its injective semilinear $j_M$ maps the target
lattice onto $L_j^{\rm can}$ by its exact generator law; fullness and
scalar extension give its rational range as precisely that $K$-span.
The theta equality then maps the rational theta submodule onto
$T_j^{\rm can}$ in both directions.  Localization commutes with this
injective map: if $s\notin\mathfrak p$ and $sy$ is in its range, divide
by the nonzero field element $s$ to lift $y$; then injectivity lifts
membership in the localized submodule.  The induced maps on each quotient
are isomorphisms, preserving its module length.  Applying this twice to
the two terms of the original signed local index proves the asserted
$I_j$ equality, with no finite-length hypothesis added.

All these definitions depend only on $p,d,W_0$.  Replacing the zero
placeholders by them changes none of the spaces, maps or sets used to
define them.  Thus one numerical record $n$ satisfies
\[
 \forall\eta\ge0\ \forall V,S,N,G,
 \quad \mathrm{NumOK}\bigl(\mathrm{assemble}
             (p.\mathrm{what}_2(d,W_0,n,\eta),V,S,N,G)\bigr).
 \tag{PKT.28}
\]
The quantifier over $\eta$ occurs after the choice of $n$.
% END PKT26

% BEGIN PKT27
\paragraph{PKT27 (full original $\mathrm{rowPk\_nonempty}$).}\label{claim:PKT27}
On the original premises stated in PKT22, including conjugation-invariance
of $W_t$, there is an actual $q:\mathrm{RowPk}(w,z_0)$.  Take its $y$
from PKT22, its $p,t$ from PKT21 applied by PKT23, its period data from
PKT24, its $W_0$, tuple and Sha clauses from PKT25, and its numerical
record from PKT26.  The tuple identity is
$y.\mathrm{tup}=\mathrm{ixOf}$ and the full matrix identity is
$p.m_L=m_{LU}$ at every tuple and every natural index, including the
zero cases.  The final two fields are universally quantified:
\[
 \forall\eta\ge0\ \forall V,S,N,G,
   \qquad \mathrm{KapDef}\ \text{and}\ \mathrm{NumOK}.
 \tag{PKT.29}
\]
PKT25 proves the first because the inserted packet preserves both weight
functions; PKT26 proves the second.  This supplies every field of the
original \texttt{RowPk}, rather than just its comparison subrecord.
The given realization $z_0$ and its original column choices are unchanged.
% END PKT27

\subsection{Copying that same package to a fitting later stage}

% BEGIN PKT28
\paragraph{PKT28 (the original realization copied with the same finite group).}\label{claim:PKT28}
Fix the $q$ of PKT27, and then fix an arbitrary $o:q.\mathrm{Over}$.
Assume exactly the original extra premises of \texttt{exists\_moved}:
$\mathrm{Holds}(z_0.R,z_0.F)$, $\mathrm{ExAt}(z_0.R)$,
$\mathrm{InGeo}(z_0.\mathrm{geo})$, $\mathrm{GrpTrue}(Z_0.N,f)$ and
$o.\mathrm{Fits}$.  The last premise supplies
$\mathrm{GrpTrue}(o.Z.N,f)$ for the same $f$.

Identify the two finite groups by the original place-label bijections.
They preserve the underlying primes, residue-field identifications and
both distinguished-prime families.  The type-idempotent coefficients agree
because both group records describe the same original $f$.  Thus the
product group isomorphism $s$ preserves the coefficient function of the
good idempotent.  Copy $z_0$ as follows: keep its actual $X_g$, trace data,
coefficient field and field embedding; copy its geometric core and its
real adelic record along the identity of the underlying fields and ranks;
compose $c_G$ with $s$; and copy its geometric half to the same original
spaces.  This is the specified $\mathrm{cpDat}$, not a new geometric
existence choice.

Here are the nontrivial field checks.  The group record uses the transported
place bijection and the same residue-field maps, so its reduction and
character laws hold pointwise.  In the good-character record, both the
character test
$\sum_g e_g\sigma(\chi(g))=\sigma(\chi(1))$ and the coefficient equality
for $\mathrm{idemOf}(X_g)$ are unchanged by $s$.  The level record keeps
its six clauses: actual level law, chain frame, idele character, group
primes, $Q$-set and base frame.  The eleven \texttt{Keep} clauses keep
the same group and two level characters, standard lattices, signs, frames,
base matrices, top pairing, nondegeneracy guard and three coefficient-field
conditions.  The weak characterization keeps the original base models,
inert piece, three full-level clauses, scalars, tuple data, phases and
all-place condition.  Thus all \texttt{Pins} hold on the copy.  The
primitive rows, auxiliary rows, finite coordinate laws and separation in
\texttt{Holds} are the same assertions about the same functions and
coefficient matrices.  The given \texttt{ExAt}, \texttt{InGeo} and
canonical-core properties similarly keep their original witnesses.
For the core laws, the multiplication, cusp multiplication, integral
$Z$-lattice, nonzero cusp and exponent clauses all use those unchanged
objects.  Therefore the copy preserves exactly the original conditional
\texttt{CoreP} implication as well.
% END PKT28

% BEGIN PKT29
\paragraph{PKT29 (all comparison maps and metrics survive the copy).}\label{claim:PKT29}
Copy $y$ with the identical tuple map and inverse and the identical frame
homomorphisms $P_Y,P_W$.  Its remaining fields are the same rank, range,
injectivity and frame equalities of PKT22.  The original groups, scalar
fields, Hermitian spaces, local representations, real actions and finite
adelic actions have not changed.  Consequently their multiplicity spaces
are equal as subspaces of the same map spaces: every transformation,
smoothness, cuspidality, growth and continuity condition has the same
underlying formula.  Their good parts are equal because the idempotent
coefficients and character laws were checked in PKT28.  Hence the original
occurrence sets, row parts and lift parts are equal, with the same direct
and conjugate type conventions.  The exact same catalogue and
representatives still satisfy \texttt{Reps}.

The row subspaces $\mathrm{EEc}$, $\mathrm{CC}$, $\mathrm{CCP}$,
$\mathrm{CCX}$ and $\mathrm{CCXP}$ are therefore equal on the two stages.
Their measures, fundamental domains, local tensors and constants $n_b,d_X$
are literally the old ones, so their pairings are equal, including any
choice-based Schur definition.  The original $\mathrm{ioC}$ and
$\mathrm{ioX}$ have the same values: their domain casts change only the
proof of membership in these equal subspaces.  Copying $j_E,j_M,j_X$
therefore preserves each injection, range, metric, adjoint transition,
vanishing, labelled-basis and coordinate field of both \texttt{ChT}s.

For the lift, $\mathrm{Wat}$, $\mathrm{MultW}$ and $\mathrm{Wc}$ have the
same defining level and good-cut conditions.  Copy $e_W,j_{M_j},j_{M_t}$
and every lattice generator.  The norm constant is still $n_0$, and
(PKT.18) is the same integral membership equivalence.  The theta clause
continues to quantify over the same $w$ with the same restriction and
image conditions; no new $w$ is selected.  The type intersections and
coordinate reading are unchanged.  Thus the full \texttt{LfT} is
preserved with the identical maps.

Finally the three-stage complex data $\mathrm{toCx3}$ are identical:
their block matrices, rational kernels, generalized eigenspaces, norms and
frames are formed from the same finite data and the same maps just copied.
The sharp dimension equality, semilinear row basis and every $X_o(v)$
in PKT20 are therefore unchanged.  This proves the full transported
\texttt{RrT}.  It also proves directly from its definition that
\[
 \eta_{\rm Of}(R',F';j,v)=\eta_{\rm Of}(R,F;j,v)
 \quad\text{for every }j,v.\tag{PKT.30}
\]
In the valid range this is the same $\eta_J$ of the same original rational
cut at the same twisted place and tuple; outside it both sides are zero.
The forced-content place set $W_1$ is likewise identical, since its tests
use those same index functions, dual indices and the same supplied error
function.  The original $D_{\rm all}$ remains the dimension of the same
full multiplicity space.  These checks give the complete transported
\texttt{PkT}, with identical \texttt{tA}, representatives, both chains'
three reading maps, and the lift's $e_W,j_{M_t}$.
% END PKT29

% BEGIN PKT30
\paragraph{PKT30 (full $\mathrm{exists\_moved}$ and the exact reader boundary).}\label{claim:PKT30}
Under PKT28's original premises there is
\[
 m:\mathrm{Moved}(o),\qquad m.\mathrm{Copy}\ \wedge\ m.\mathrm{Coh}.
 \tag{PKT.31}
\]
Take $m.z$ to be PKT28's copied realization.  Its intersection/holds,
$\mathrm{ExAt}$ and $\mathrm{InGeo}$ fields are the given properties
transported there.  Its $X_g$ and trace fields are unchanged.  Its core
field is exactly the conditional core implication proved in PKT28.
Its $\eta$ and $W_1$ fields are PKT29.  Its $y$, tuple identity and
comparison packet are the copied ones in PKT29, and its \texttt{tA}
field records equality of the original complexification maps.  This is
every field of \texttt{Moved}.

The two \texttt{Copy} clauses say that its geometric core and real record
are the prescribed copies of $z_0$; they hold by the construction.  The
nine \texttt{Coh} clauses assert identity of the tuple inverse, row frame,
all lift frames, representatives, both chains' $j_E,j_M,j_X$, and the
lift's $e_W,j_{M_t}$.  They hold simultaneously because PKT29 copied
those functions, rather than choosing new coordinates or orthonormal
bases.  This proves the full original endpoint for every later $o$
satisfying its original $\mathrm{Fits}$ guard.

RDR12 and RDR20--RDR21 may now be applied to this same $m$ and $o$ to
construct the actual level-zero rational side and \texttt{ReadPk};
RPI15 and DI supply the already reconstructed first division comparison
on the same underlying objects.  These joins add no row-integrability
premise.  Full \texttt{ReadL} still requires precisely
$o.\eta(j,v)=q.\eta(j,v)$ for $1\le j\le m$: (PKT.30) identifies the
copied realized parameter with $q.\eta$, not an arbitrarily supplied
$o.\eta$ with that realized parameter.  The outer common-family production
of $z_0$, its original analytic/geometric \texttt{Inputs/Pins/Holds},
$\mathrm{ExAt}/\mathrm{InGeo}$, the fitting division/rational tests,
and the stronger second division-row theta comparison remain independent
obligations wherever they have not already been supplied.  None is
inferred from (PKT.31).
% END PKT30
